\documentclass[11pt]{amsart}

\usepackage[T1]{fontenc}
\usepackage{lmodern}
\usepackage{amsmath,amssymb,amsthm}
\usepackage{mathtools}
\usepackage{stmaryrd}      % \llbracket \rrbracket
\usepackage{pifont}        % \ding
\usepackage[colorlinks=true,linkcolor=blue,citecolor=blue,urlcolor=blue]{hyperref}
\usepackage[protrusion=true,expansion=false]{microtype}

% ---- theorem environments -------------------------------------------------
\theoremstyle{plain}
\newtheorem{theorem}{Theorem}[section]
\newtheorem{proposition}[theorem]{Proposition}
\newtheorem{lemma}[theorem]{Lemma}
\newtheorem{corollary}[theorem]{Corollary}
\theoremstyle{definition}
\newtheorem{definition}[theorem]{Definition}
\newtheorem{example}[theorem]{Example}
\theoremstyle{remark}
\newtheorem{remark}[theorem]{Remark}

% ---- notation -------------------------------------------------------------
\newcommand{\Set}{\mathsf{Set}}
\newcommand{\Cont}{\mathsf{Cont}}
\newcommand{\Kl}{\mathrm{Kl}}
\newcommand{\coKl}{\mathrm{coKl}}
\newcommand{\Alg}{\mathrm{Alg}}
\newcommand{\Coalg}{\mathrm{Coalg}}
\newcommand{\op}{\mathrm{op}}
\newcommand{\id}{\mathrm{id}}
\newcommand{\str}{\mathrm{st}}
\newcommand{\lv}{\mathrm{lv}}
\newcommand{\arr}{\mathrm{arr}}
\newcommand{\first}{\mathrm{first}}
\newcommand{\Ext}[1]{\llbracket #1\rrbracket}
\newcommand{\cont}[2]{(#1,#2)}
\newcommand{\Maybe}{\mathsf{Maybe}}
\newcommand{\Pf}{\mathsf{Pf}}
\newcommand{\Writer}{\mathsf{Writer}}
\newcommand{\List}{\mathsf{List}}
\newcommand{\Arr}{\mathsf{Arr}}
\newcommand{\rcomp}{\mathbin{>\!\!>\!\!>}}
\DeclareMathOperator{\Nat}{Nat}
\DeclareMathOperator{\Hom}{Hom}

\begin{document}

\title[Effects and coeffects entwine on containers]{Effects and coeffects entwine on containers:\\
one monad, two feeds, and the branching obstruction}

\author{MacBeth}
\address{Kodamai}
\email{}

\date{31 July 2026}

\subjclass[2020]{18C15, 18M05, 18C50, 68Q55}
\keywords{container, polynomial functor, distributive law, entwining, monad, comonad,
effect, coeffect, arrow, Freyd category, bialgebraic semantics}

\begin{abstract}
A single monad $M$ on $\Set$ acts on the category $\Cont$ of containers along either
axis: on \emph{shapes} it gives the Ahman--Bauer effect monad $T_M$; on \emph{positions}
it gives, by contravariance, a coeffect comonad $G_M$. We ask how these two feeds of one
monad interact. They always entwine: for every $M$ with a $\prod$-cointerpretation lifting
there is a canonical mixed distributive law $\lambda\colon T_M G_M \Rightarrow G_M T_M$
--- the oplax product-comparison of $M$, read on positions --- making $G_M$ a comonad on
$T_M$-algebras and $T_M$ a monad on $G_M$-coalgebras. This Turi--Plotkin bialgebra holds
unconditionally, nondeterminism included. The reverse orientation is another story: the
compositor $\kappa\colon G_M T_M \Rightarrow T_M G_M$ that would let one \emph{sequentially
compose} effect--coeffect programs $G_M p \to T_M q$ into a category exists if and only if
$M$ is \emph{non-branching} (every operation has arity at most one), a class we identify
exactly as the writer-with-absorbing-exceptions monads $MX \cong E + A\times X$. There the
arrows form a Hughes arrow, equivalently a Freyd category over $(\Cont,\times)$; for
branching $M$ they fail, through two independent axioms. Branching obstructs the
\emph{orientation} of the effect--coeffect interaction, not the interaction itself.
\end{abstract}

\maketitle

%============================================================================
\section{Introduction}
%============================================================================

A monad is a model of an \emph{effect}: a way of building a computation that returns a
value while also doing something --- failing, choosing, logging, diverging. A comonad is
a model of a \emph{coeffect}: a way of consuming a computation that reads a value together
with some ambient \emph{context} --- an environment, a history, a resource budget. The two
notions are formally dual, and a long line of work in programming semantics has asked how
they interact when both are present at once: how does one \emph{compose} a program that
reads context and produces effects with another of the same shape? The classical answers
--- Hughes arrows \cite{Hughes00}, Freyd categories \cite{PowerRobinson97,Atkey11},
biKleisli categories of a monad and a comonad with a distributive law
\cite{PowerWatanabe02,UustaluVene08}, and the bialgebraic operational semantics of Turi and
Plotkin \cite{TuriPlotkin97} --- all live over an arbitrary base category and take the
monad and the comonad as \emph{independent} data, related (if at all) by a distributive law
that must be supplied by hand.

This paper studies the case where the monad and the comonad are not independent but are the
\emph{two feeds of a single monad}, and where the base category is the category $\Cont$ of
containers (equivalently, polynomial functors in one variable). The setting makes the
interaction canonical, and canonicity turns out to have sharp consequences.

\paragraph{The two feeds.}
A container $\cont S P$ is a set $S$ of \emph{shapes} together with a family $P\colon S\to
\Set$ of \emph{positions}; it presents the polynomial endofunctor $\Ext{S,P}Y = \sum_{s\in
S} Y^{P(s)}$ on $\Set$. A morphism of containers is a map \emph{forward} on shapes and
\emph{backward} on positions --- positions are contravariant --- and this single fact about
variance is the hinge of the whole paper. Given any monad $M$ on $\Set$ (with a mild
support hypothesis, made precise in \S\ref{sec:prelim}), there are exactly two ways to feed
$M$ into a container, and the variance decides the outcome:
\begin{itemize}
  \item \emph{Along the shapes}, $T_M\cont S P = \cont{MS}{P^\star}$ applies $M$ to the
  shape set and extends the positions to match (Ahman--Bauer \cite{AhmanBauer24}, Thm.~6.3).
  Shapes are covariant, so $T_M$ is a \textbf{monad} on $\Cont$ --- the \emph{effect} feed.
  \item \emph{Along the positions}, $G_M\cont S P = \cont S{M\circ P}$ applies $M$ to each
  position fibre and leaves the shapes fixed (the \emph{transfer}). Positions are
  contravariant, so $G_M$ is a \textbf{comonad} on $\Cont$ --- the \emph{coeffect} feed.
\end{itemize}
There is no third axis and no freedom in the outcome. The reader should hold onto the
asymmetry: it is the same variance --- positions run backward --- that makes $G_M$ a comonad
and, we will find, forces the single direction in which the effect feed can pass through the
coeffect feed.

\paragraph{The question, and the answer.}
Do the two feeds $T_M$ and $G_M$ of one monad $M$ interact via a distributive law, and if
so, what does that law model? We give a complete answer, and it is a \emph{dichotomy of two
dual laws} that coincide only in a degenerate regime.

\begin{theorem}[Main dichotomy]\label{thm:main}
Let $M$ be a $\Set$-monad with a $\prod$-cointerpretation predicate lifting (for instance
$\Pf$, $\Maybe$, $\Writer$, or $\mathrm{Id}$), with effect monad $T_M$ and coeffect comonad
$G_M$ on $\Cont$ as above. Then:
\begin{enumerate}
  \item \emph{(Bialgebra face, all $M$.)} There is a canonical mixed distributive law
  $\lambda\colon T_M G_M \Rightarrow G_M T_M$, the oplax product-comparison of $M$ read on
  positions, satisfying the four entwining axioms for \emph{every} such $M$. Consequently
  $G_M$ lifts to a comonad on $T_M$-algebras and $T_M$ lifts to a monad on
  $G_M$-coalgebras: a Turi--Plotkin-style $\lambda$-bialgebra.
  \item \emph{(Arrow face, non-branching $M$ only.)} The reverse compositor
  $\kappa\colon G_M T_M \Rightarrow T_M G_M$ satisfies the four axioms --- equivalently, the
  effect--coeffect programs $p\rightsquigarrow q := \Cont(G_M p,\,T_M q)$ form a category,
  indeed a Hughes arrow / Freyd category over $(\Cont,\times)$ --- \emph{if and only if $M$
  is non-branching}, meaning every shape of $M$ has arity at most one. The sole obstructed
  axiom is associativity.
  \item \emph{(Classification of the positive class.)} A cartesian $M$ is non-branching if
  and only if it is a \emph{writer-with-absorbing-exceptions} monad
  \[
    MX \;\cong\; E + A\times X,\qquad (A,\cdot,e)\ \text{a monoid},\quad
    (E,\odot)\ \text{a left $A$-set},
  \]
  with writer unit $x\mapsto(e,x)$ and the log $A$ acting on the exceptions $E$; the arrow
  category exists throughout this class.
\end{enumerate}
\end{theorem}

Parts (1)--(3) are Theorems~\ref{thm:entwine}, \ref{thm:arrowsA}--\ref{thm:arrowsB}, and
\ref{thm:classify} below. The single monad $M$ carries two comparison maps between its
composites $T_M G_M$ and $G_M T_M$, and they behave oppositely: the oplax one runs
$T_M G_M \Rightarrow G_M T_M$ and is always a law; the lax one runs the other way and is a
law only when $M$ does not branch. \emph{Same monad, same two feeds, same category $\Cont$:
the unification of effects and coeffects exists unconditionally as a bialgebra, and its
categorical packaging as an arrow calculus is exactly what branching obstructs.} The
obstruction is not to the interaction but to its \emph{orientation}.

\paragraph{Why this is worth stating.}
Three points give the result its edge.

First, the bialgebra face \emph{keeps the branching monads}. A recurring worry about
effect--coeffect calculi is that the interesting computational monads --- nondeterminism
$\Pf$, probability, lists --- are precisely the ones that break the sequential composition.
They do break the \emph{arrow}, but they do not break the \emph{bialgebra}: part (1) holds
for $\Pf$ and $\List$ exactly as for $\Maybe$. So the denotational, bialgebraic unification
of effects and coeffects on containers is available for every $M$, and the affirmative
answer to ``is there a Turi--Plotkin bialgebra here?'' is unconditional.

Second, the obstruction is \emph{named}. Part (3) does not merely say ``non-branching''; it
identifies the class as the writer-with-absorbing-exceptions monads $E + A\times(-)$, a
recognisable and closed family (the writer transformer over the exception monad, twisted by
an action). ``Non-branching'' is moreover a genuinely new dividing line, provably orthogonal to
commutativity and affineness (Proposition~\ref{prop:independence}): $\Pf$ is commutative and branching,
$E+(-)$ with $|E|>1$ is non-commutative and non-branching, the distribution monad is affine
and branching.

Third, branching disables the arrow \emph{twice over}, through two independent axioms. The
compositor $\kappa$ fails associativity (Theorem~\ref{thm:arrowsA}) through the
multiplication of $M$ \emph{merging} leaves (a \emph{union of products versus product of
unions} discrepancy); and, separately, the effect monad $T_M$ has no \emph{natural} strength
for the product on $\Cont$ (Lemma~\ref{lem:strong}), failing through leaf-\emph{symmetry}
rather than merging. Both failures are equivalent to the arity condition, by different
arguments --- a pleasingly redundant obstruction.

\paragraph{Method.}
Because both feeds are identity on shapes over $MS$ (for $T_M$) or on $S$ (for $G_M$), every
law and every axiom localises to a \emph{position-level equation in $\Set$}, with backward
maps composing contravariantly. At a shape with leaves $b$ and position sets $Z_b$, the two
composites $T_M G_M$ and $G_M T_M$ read as $M(\prod_b Z_b)$ and $\prod_b M(Z_b)$, and the
only canonical map between them is the oplax product-comparison
$\str\colon M(\prod_b Z_b)\to\prod_b M(Z_b)$. The entwining axioms are then naturality
squares for $M$'s unit and multiplication; the arrow's obstruction is that the \emph{reverse}
(lax) comparison fails to commute with the leaf-merging in $\mu^M$; and the strength
obstruction is a Yoneda computation showing a natural distributivity map must read a single
distinguished leaf, which leaf-symmetry then forbids. The classification is a
polynomial-functor argument (arity $\le 1$ is degree $\le 1$) followed by a monad-structure
pin.

\paragraph{Outline.}
Section~\ref{sec:prelim} fixes containers, the two liftings, and the entwining vocabulary.
Section~\ref{sec:feeds} presents the effect monad $T_M$ and the coeffect comonad $G_M$ and
the sense in which they are one monad's two feeds. Section~\ref{sec:bialgebra} proves the
bialgebra face --- the law $\lambda$ for all $M$ --- and Section~\ref{sec:onedirection}
shows the reverse orientation fails, isolating branching as the obstruction.
Section~\ref{sec:arrow} builds the arrow face: the biKleisli category
(Theorem~\ref{thm:arrowsA}), the Freyd/Hughes structure and the strength obstruction
(Theorem~\ref{thm:arrowsB}). Section~\ref{sec:classify} classifies the positive class as
writer-with-absorbing-exceptions and closes the associativity axiom across it
(Theorem~\ref{thm:classify}). Section~\ref{sec:lean} reports the machine-checked instance.
Section~\ref{sec:related} places the result among its neighbours, and
Section~\ref{sec:conclusion} concludes.

\paragraph{Provenance.}
The mathematical results collected here were established in a sequence of the author's proof
notes; the effect monad $T_M$ is due to Ahman and Bauer \cite{AhmanBauer24}, and the
entwining and arrow formalism is classical (Beck \cite{Beck69}; Power--Watanabe
\cite{PowerWatanabe02}; Brzezi\'nski--Majid \cite{BrzezinskiMajid98}; Uustalu--Vene
\cite{UustaluVene08}). What is new is the identification of the law welding one monad's two
container feeds, its forced orientation, the branching obstruction, and the classification
of the positive class. Attribution is given in place. Two axioms are proved in coordinates
and machine-verified on the flagship monads but not written out as a fully symbolic chase over
an arbitrary predicate lifting; these mechanical gaps are flagged honestly where they occur
(Remarks~\ref{rem:e2gap} and \ref{rem:scope}).

%============================================================================
\section{Preliminaries}\label{sec:prelim}
%============================================================================

We work with containers and with $\Set$-monads carrying a support structure. This section
fixes notation and recalls the entwining vocabulary; a reader fluent in polynomial functors
and distributive laws can skim it.

\subsection{Containers}
A \emph{container} is a pair $\cont S P$ with $S\in\Set$ (\emph{shapes}) and $P\colon S\to
\Set$ (\emph{positions}). A \emph{morphism} $\cont S P \to \cont T Q$ is a pair $(u,f)$ with
$u\colon S\to T$ \emph{forward} and, for each $s\in S$, a map $f_s\colon Q(u\,s)\to P(s)$
\emph{backward}; composition is forward on shapes and backward (in the opposite order) on
positions. Containers and their morphisms form a category $\Cont$. The \emph{extension}
functor
\[
  \Ext{-}\colon \Cont \longrightarrow [\Set,\Set],\qquad
  \Ext{S,P}Y = \sum_{s\in S} Y^{P(s)},
\]
is fully faithful, with image the \emph{polynomial} endofunctors; we freely identify a
container with the functor it presents. Finite products in $\Cont$ are
\[
  \cont S P \times \cont T Q = \bigl(S\times T,\ (s,t)\mapsto P(s)\sqcup Q(t)\bigr),
\]
matching $\Ext{p\times c}Y = \Ext p Y\times\Ext c Y$; this is the cartesian tensor used in
\S\ref{sec:arrow}.

\subsection{Set-monads with support and the predicate lifting}\label{sec:prelim-lifting}
Let $M=(M,\eta,\mu)$ be a monad on $\Set$ presented as a polynomial functor
\[
  M X \;=\; \sum_{s\in M1} X^{\,\mathrm{ar}(s)},\qquad \mathrm{ar}\colon M1\to\Set,
\]
so each element $m\in MX$ has a \emph{support} of \emph{leaves} $b\in\lv(m)$ carrying labels
$x_b\in X$. We say $M$ is \emph{non-branching} if $|\mathrm{ar}(s)|\le 1$ for every shape
$s$ (equivalently every $m$ has $|\lv(m)|\le 1$), and \emph{branching} otherwise. Flagship
examples: $\Maybe = 1+(-)$, exception $E+(-)$, $\Writer_N = N\times(-)$, and $\mathrm{Id}$
are non-branching; finite powerset $\Pf$ and $\List$ are branching.

To lift $M$ into positions we need a \emph{predicate lifting} $(-)^\star$ turning a family
$P\colon S\to\Set$ into $P^\star\colon MS\to\Set$. We use the \emph{$\prod$-cointerpretation}
of Ahman and Bauer \cite[\S6.1,\S6.5]{AhmanBauer24}:
\[
  P^\star(m) \;=\; \prod_{b\in\lv(m)} P(x_b),
\]
the universal (product-based) lifting. Its unit datum is the projection $i$ out of the
singleton product $P^\star(\eta_S s)=\prod_{\{*\}}P s = Ps$, and its multiplication datum is
the restriction $j$ that reindexes a product along the leaf-map of $\mu^M$. The monads for
which this weak Mendler algebra is defined include all the flagship examples above; the
reader monad $A^K$ is \emph{not} one, which is why the lifting is taken weak. We call such
$M$ a \emph{$\prod$-cointerpretation monad}; this is the standing hypothesis throughout, and
\emph{cartesian} means $\mu^M$ (and $\eta^M$) are cartesian natural transformations
(preserve the leaf-count).

\begin{remark}[``Non-branching'' is not Kock-affine]\label{rem:affine-clash}
Non-branching means \emph{arity $\le 1$}, i.e.\ degree $\le 1$ as a polynomial (a constant
term plus a linear term). It is \emph{not} affineness in Kock's sense $M1\cong 1$: here the
shape set $M1$ may be arbitrarily large. The two coincide only trivially ($M1\cong1$ forces
$M=\mathrm{Id}$). We say \emph{non-branching} or \emph{arity-$\le 1$} throughout and never
\emph{affine}, to avoid the clash.
\end{remark}

\subsection{Mixed distributive laws and entwinings}\label{sec:prelim-entwine}
Let $T$ be a monad and $G$ a comonad on a category $\mathcal C$. A \emph{mixed distributive
law} (or \emph{entwining}) of $G$ over $T$ is a natural transformation
$\lambda\colon TG\Rightarrow GT$ satisfying four axioms --- two expressing compatibility with
the unit and multiplication of $T$, two with the counit and comultiplication of $G$
(Beck \cite{Beck69}; Power--Watanabe \cite{PowerWatanabe02}; the entwining-structure
language is Brzezi\'nski--Majid \cite{BrzezinskiMajid98}). Such a $\lambda$ is exactly the
datum lifting $G$ to a comonad on the category $\Alg(T)$ of $T$-algebras and $T$ to a monad
on the category $\Coalg(G)$ of $G$-coalgebras; a compatible $T$-algebra-and-$G$-coalgebra is
a \emph{$\lambda$-bialgebra} in the sense of Turi and Plotkin \cite{TuriPlotkin97}. We write
the four axioms E1--E4 (unit-$T$, mult-$T$, counit-$G$, comult-$G$) and, for the opposite
orientation $\kappa\colon GT\Rightarrow TG$, their mirror E1$'$--E4$'$; here E2$'$ is the
mult-$T$/associativity axiom.

Dually, a mixed distributive law $\kappa\colon GT\Rightarrow TG$ lifts $G$ to a comonad on
the Kleisli category $\Kl(T)$, whose \emph{coKleisli} category is the \emph{biKleisli}
category of the pair: its objects are those of $\mathcal C$, its morphisms $p\to q$ are the
maps $Gp\to Tq$, and composition uses $\kappa$ to thread the effect of one map past the
coeffect of the next \cite{PowerWatanabe02,UustaluVene08}. When $T$ is strong and $G$
costrong compatibly with $\kappa$, this biKleisli category is a \emph{Hughes arrow}
\cite{Hughes00}, equivalently a \emph{Freyd category} over the cartesian base
\cite{PowerRobinson97,Atkey11,JacobsHeunenHasuo09}. These are the two structures the paper
produces from the two feeds of one $M$.

%============================================================================
\section{The two feeds of one monad}\label{sec:feeds}
%============================================================================

We recall the two liftings and the exact sense in which they are dual halves of one monad.

\subsection{The effect monad}
Ahman and Bauer \cite[Thm.~6.3]{AhmanBauer24} show that a $\prod$-cointerpretation monad $M$
lifts along shapes to a monad on $\Cont$:
\[
  T_M\cont S P = \cont{MS}{P^\star},\qquad P^\star(m)=\prod_{b\in\lv(m)}P(x_b).
\]
The unit $\eta^T$ covers $\eta^M$ on shapes and is backward the singleton projection $i$; the
multiplication $\mu^T$ covers $\mu^M$ and is backward the restriction $j$. Because shapes are
covariant, $T_M$ is a monad. On extensions $\Ext{T_M(S,P)}Y = \sum_{m\in MS} Y^{P^\star(m)}$:
$T_M$ is the container of \emph{$M$-structured computations over the operations of $\cont S
P$}, with positions the tuples of operands, one per leaf.

\subsection{The coeffect comonad}
Dually, apply $M$ to the position fibres:
\[
  G_M\cont S P = \cont S{M\circ P}.
\]
The counit $\varepsilon$ is backward $\eta^M$ and the comultiplication $\delta$ is backward
$\mu^M$; shapes are untouched. Because positions are \emph{contravariant}, the covariant
monad $M$ becomes a \emph{comonad} $G_M$ on $\Cont$. On extensions $\Ext{G_M(S,P)}Y =
\sum_{s} Y^{M(P s)}$: an operation of arity $P(s)$ acquires an $M$-structured arity, a
\emph{context} wrapped around every operand slot. (For $M=\Maybe$, $G_M$ adjoins a fresh
basepoint to each fibre --- a default or exception argument on every operation.)

The two liftings are welded together by the fibrewise $(-)^\op$ that makes a container a
container: apply $M$ where the variance runs forward and monads stay monads; apply it where
the variance runs backward and monads become comonads. There is no third axis, and no choice
in the outcome. The rest of the paper studies how these two feeds of \emph{the same} $M$
interact.

%============================================================================
\section{The bialgebra face: the two feeds always entwine}\label{sec:bialgebra}
%============================================================================

This is the headline. We show the two feeds interact via a canonical mixed distributive law
$\lambda\colon T_M G_M\Rightarrow G_M T_M$, valid for \emph{every} $\prod$-cointerpretation
$M$, branching or not, and identify what it models.

\subsection{The two composites and the canonical comparison}
Both $T_M G_M$ and $G_M T_M$ are the identity on shapes over $MS$, so it suffices to read
their positions at a fixed shape $m\in MS$. Localise at $m$, write $Z_b := P(x_b)$ for the
position set at leaf $b\in\lv(m)$, and compute:
\[
  \renewcommand{\arraystretch}{1.4}
  \begin{array}{c|c}
    \text{functor} & \text{positions at } m\\\hline
    G_M\,T_M & M\bigl(P^\star m\bigr)=M\bigl(\textstyle\prod_b Z_b\bigr)\\
    T_M\,G_M & (M\circ P)^\star(m)=\textstyle\prod_b M(Z_b)
  \end{array}
\]
($G_M T_M$: run $T_M$, then fatten its positions with $M$. $T_M G_M$: fatten first, then
take the product-lifting of the $M$-fattened positions.) Between $M(\prod_b Z_b)$ and
$\prod_b M(Z_b)$ there is exactly one canonical map, the \emph{oplax product-comparison}
\[
  \str_{(Z_b)}\colon\ M\bigl(\textstyle\prod_b Z_b\bigr)\longrightarrow\prod_b M(Z_b),
  \qquad w\longmapsto\bigl(M(\pi_b)\,w\bigr)_b,
\]
built from $M$ and the product projections alone. It exists for \emph{every} functor $M$ by
the universal property of the target product --- it is the structure map exhibiting $M$ as an
oplax monoidal functor for $(\Set,\times)$ --- and no algebra structure on $M$ is used to
write it down.

\paragraph{The direction is forced.}
$\str$ runs $G_M T_M$-positions to $T_M G_M$-positions. Container morphisms run their
position map \emph{backward}, so $\str$ is the backward part of a morphism the \emph{other}
way,
\[
  \lambda_X\colon\ T_M G_M X\longrightarrow G_M T_M X,\qquad
  \text{forward }\id_{MS},\quad\text{backward }\str,
\]
a $2$-cell $\lambda\colon T_M G_M\Rightarrow G_M T_M$: the standard orientation of a mixed
distributive law of the comonad $G_M$ over the monad $T_M$. The naive opposite orientation
would need the reverse map $\prod_b M(Z_b)\to M(\prod_b Z_b)$, which is not canonical and, we
show in \S\ref{sec:onedirection}, does not assemble into a law.

\subsection{The entwining theorem}

\begin{theorem}[The two feeds entwine]\label{thm:entwine}
Let $M$ be a $\prod$-cointerpretation $\Set$-monad, $T_M,G_M$ its two container liftings, and
$\lambda\colon T_M G_M\Rightarrow G_M T_M$ the $2$-cell backward $\str$. Then $(T_M,G_M,
\lambda)$ is an entwining structure on $\Cont$: $\lambda$ is natural and satisfies the four
entwining axioms. Consequently $G_M$ lifts to a comonad on $\Alg(T_M)$ and $T_M$ lifts to a
monad on $\Coalg(G_M)$, exhibiting a $\lambda$-bialgebra in the sense of Turi and Plotkin.
\end{theorem}

\begin{proof}
Everything is identity on shapes, so each axiom localises to a position-level equation in
$\Set$, with backward maps composing contravariantly; localise at $Z_b$ and write
$\str=(M\pi_b)_b$.
\begin{itemize}
  \item \textbf{Counit-$G$} ($\varepsilon_G$ against $\lambda$). Backward,
  $M(\pi_b)\circ\eta_{\prod Z}=\eta_{Z_b}\circ\pi_b$: the naturality of $\eta$ at $\pi_b$.
  \item \textbf{Unit-$T$}. At a single shape the product collapses to one factor; there
  $\str=M(\id)=\id$ and the Mendler $i$ is the singleton projection $=\id$, so both sides are
  $\id_{MZ}$.
  \item \textbf{Comult-$G$} ($\delta_G$ against $\lambda$). Backward,
  $M(\pi_b)\circ\mu_{\prod Z}=\mu_{Z_b}\circ MM(\pi_b)$: the naturality of $\mu$ at $\pi_b$.
  Threading the two paths componentwise, both equal $(\mu_{Z_b}\circ MM(\pi_b))_b$.
  \item \textbf{Mult-$T$}. The only axiom touching the Mendler $j$. Both paths are built from
  $\str$ and the reindexing of a product along the leaf-map of $\mu^M$; since $\str$ is
  componentwise $M$ applied to a projection, it commutes with that reindexing
  ($M(\pi_{\rho(b,c)})=M(\pi_{(b,c)})\circ M(\text{reindex})$), and the two backward maps
  coincide. This is the naturality of $\str$ with respect to product-reindexing, i.e.\ the
  $j$-naturality diagram of Ahman--Bauer \cite[Def.~6.2]{AhmanBauer24}.
\end{itemize}
Naturality of $\lambda$ in $X$ is naturality of $\str$ and of the projections. The lifting to
$\Alg(T_M)$ and $\Coalg(G_M)$ is the standard consequence of the entwining axioms
\cite{PowerWatanabe02}.
\end{proof}

\begin{remark}[The one-line meaning of $\lambda$]
The entwining says one thing: \emph{$M$ oplax-preserves products, coherently with its unit
and multiplication, transported onto positions.} Fibrationally, $G_M=(M^\op)_*$ is vertical
and $T_M$ covers the base monad $M$; $\lambda$ is the Beck--Chevalley-type $2$-cell comparing
\emph{apply $M$ on the base, then fatten positions} against \emph{fatten positions, then apply
$M$ on the base}. The comparison exists one way for free --- that is $\str$ --- and the four
axioms are exactly its coherence with the two monad structures $M$ carries on the two feeds.
\end{remark}

\begin{remark}[The mult-$T$ index-chase]\label{rem:e2gap}
Axiom mult-$T$ is the only one whose full symbolic verification depends on the specific weak
Mendler $j$. The conceptual mechanism above --- $\str$ commutes with the product-reindexing
that is $j$ --- is complete; the fully spelt-out index-chase over an arbitrary $j$ is
mechanical and is not written out. It is machine-verified for the $\prod$-cointerpretation
examples across three containers, including the \emph{branching} commutative case $M=\Pf$
(union/restriction $j$), $M=\Maybe$, and $M=\Writer$ over both an abelian and a
non-commutative log. The other three axioms are complete as stated.
\end{remark}

The bialgebra face holds for \emph{every} $M$. In particular nondeterminism ($\Pf$),
lists, and --- for any $\prod$-cointerpretation presentation --- probabilistic branching all
carry the $\lambda$-bialgebra: the two feeds of a branching monad still entwine in the
standard orientation, and $G_M$ still acts on $T_M$-algebras. This is the sense in which the
denotational unification of effects and coeffects on containers is unconditional. What
branching costs is the \emph{other} orientation, to which we now turn.

%============================================================================
\section{One direction only: the branching obstruction}\label{sec:onedirection}
%============================================================================

The opposite orientation $G_M T_M\Rightarrow T_M G_M$ needs the \emph{lax} comparison
$\prod_b M(Z_b)\to M(\prod_b Z_b)$, which not every monad has; a commutative monad supplies
the cartesian one (for $\Pf$, $(S_b)_b\mapsto\prod_b S_b$, the set of choice-tuples). Even
when it exists, it fails to be a distributive law the moment $M$ branches.

\begin{example}[$\Pf$ breaks the reverse orientation]\label{ex:pf-breaks}
Take $M=\Pf$ and the container $X=\bigl(\{a,b\},\ a\mapsto\{0,1\},\ b\mapsto\{0\}\bigr)$. The
lax $G_M T_M\Rightarrow T_M G_M$ map satisfies unit-$T$, counit-$G$, and comult-$G$, but
fails mult-$T$. At the double shape $\{\{a,b\},\{a\}\}\in\Pf\,\Pf\,S$ the two paths return, on
one side, the \emph{correlated} two-element set $\{((1,0),(1)),\,((0,0),(0))\}$ --- a union of
products --- and on the other the full four-element \emph{product of unions}. The
multiplication $\mu=\bigcup$ merges the leaf $a$ shared by $\{a,b\}$ and $\{a\}$, and the
cartesian lax map cannot see that the two choices at that shared leaf must agree:
\emph{union of products $\neq$ product of unions}.
\end{example}

The obstruction is \textbf{branching, not non-commutativity}. $\Pf$ is commutative and still
fails; what breaks the law is $\mu$ identifying leaves, which needs a shape with at least two
positions so that $\bigcup$ can overlap. When $M$ has arity at most one no two leaves are
ever shared, $\mu$ can identify none, mult-$T$ holds trivially, and \emph{both} orientations
are entwinings at once. The dichotomy is clean:
\begin{center}\small
  \renewcommand{\arraystretch}{1.4}\setlength{\arraycolsep}{4pt}
  $\begin{array}{c|c|c}
    M & T_M G_M\Rightarrow G_M T_M\ (\text{oplax }\str) & G_M T_M\Rightarrow T_M G_M\ (\text{lax})\\\hline
    \text{arity}\le 1\ (\Maybe,\Writer,\mathrm{Id}) & \text{entwining }\checkmark
      & \text{entwining }\checkmark\\
    \text{branching}\ (\Pf,\List,\dots) & \text{entwining }\checkmark\ (\textbf{universal})
      & \text{fails mult-}T\ \text{\ding{55}}
  \end{array}$
\end{center}
So the interaction of the two feeds is real, and it lives in the standard orientation, where
it holds for all $M$; the orientation one might guess first is the one branching forbids. The
same variance that forced $G_M$ to be a comonad also forces the single direction in which the
monad feed can pass through it.

%============================================================================
\section{The arrow face: composing effect--coeffect programs}\label{sec:arrow}
%============================================================================

We now read the reverse orientation as a \emph{composition rule}. An \emph{effect--coeffect
arrow} $p\rightsquigarrow q$ is a container morphism $G_M p\to T_M q$: it reads its input
through the coeffect $G_M$ and delivers its output through the effect $T_M$. This is the
container analogue of a coKleisli-of-$G$-then-Kleisli-of-$T$ program $GA\to TB$, and it is
what one \emph{sequentially composes}.

\subsection{The biKleisli category and its compositor}
The arrows $G_M p\to T_M q$ are the morphisms of the coKleisli category of the comonad $G_M$
lifted to the Kleisli category $\Kl(T_M)$ --- the biKleisli category of the pair
\cite{PowerWatanabe02,UustaluVene08}. Explicitly, the identity $p\rightsquigarrow p$ is
$\eta^T_p\circ\varepsilon_p$, and the composite of $f\colon G_M p\to T_M q$ and
$g\colon G_M q\to T_M r$ is
\[
  G_M p \xrightarrow{\ \delta_p\ } G_M G_M p \xrightarrow{\ G_M f\ } G_M T_M q
      \xrightarrow{\ \kappa_q\ } T_M G_M q \xrightarrow{\ T_M g\ } T_M T_M r
      \xrightarrow{\ \mu^T_r\ } T_M r.
\]
The middle step is a natural transformation $\kappa\colon G_M T_M\Rightarrow T_M G_M$: the
composite \emph{commutes the effect $T_M$ out through the coeffect $G_M$}, from $G_M T_M q$ to
$T_M G_M q$. This is the \emph{reverse} of the proved $\lambda$ of Theorem~\ref{thm:entwine};
on positions its backward map is the \emph{lax} product-comparison $\prod_b M(Z_b)\to
M(\prod_b Z_b)$ of \S\ref{sec:onedirection}. Composition therefore requires exactly the
orientation that branching obstructs.

\begin{theorem}[Existence of the arrow category]\label{thm:arrowsA}
For a $\prod$-cointerpretation $\Set$-monad $M$ the following are equivalent.
\begin{enumerate}
  \item The effect--coeffect arrows $G_M p\to T_M q$ form a category with identity
  $\eta^T\circ\varepsilon$ and the biKleisli composition above.
  \item There is a natural $\kappa\colon G_M T_M\Rightarrow T_M G_M$ satisfying the four
  mixed-distributive-law axioms E1$'$--E4$'$.
  \item $M$ is non-branching.
\end{enumerate}
The unit laws correspond to E1$'$, E3$'$; associativity corresponds to \textbf{E2$'$} (the
mult-$T$ axiom), which is the sole axiom that branching destroys.
\end{theorem}

\begin{proof}
$(1)\Leftrightarrow(2)$ is the mixed analogue of Beck's theorem \cite{PowerWatanabe02}: a
lifting of $G_M$ to $\Kl(T_M)$ is the same datum as a $\kappa$ obeying E1$'$--E4$'$, and the
coKleisli category of the lifted comonad is (1), with the hom-sets $\Cont(G_M p,T_M q)$,
identity $\eta^T\circ\varepsilon$, and the displayed composite; the unit laws unwind to
E1$'$/E3$'$ and associativity to E2$'$, which is precisely the statement that the lifted
comonad preserves Kleisli composition. $(2)\Leftrightarrow(3)$ is the dichotomy of
\S\ref{sec:onedirection}: the lax $\kappa$ satisfies E1$'$, E3$'$, E4$'$ for every $M$ and
satisfies E2$'$ iff $M$ is non-branching, the branching failure being the union-of-products
versus product-of-unions discrepancy of Example~\ref{ex:pf-breaks}.
\end{proof}

The equivalence $(1)\Leftrightarrow(3)$ is what makes the result sharp: the existence of an
effect--coeffect arrow calculus on containers is a property of the effect monad $M$, namely
non-branching, and the obstruction is exactly the entwining axiom that branching destroys.
The signature is telling. For $M=\Pf$ the unit laws still hold and \emph{only} associativity
breaks; an explicit non-associative triple $f,g,h\colon G_M(A_1)\to T_M(A_1)$ (with
$A_1=(\{a,b\};a\mapsto2,b\mapsto1)$, all forward $a\mapsto\{a\},b\mapsto\{a,b\}$) differs at
shape $b$, position $(1,0)$, returning $\varnothing$ on one association and $\{0\}$ on the
other.

\subsection{The Freyd/Hughes structure, and the strength obstruction}
A category is not yet an arrow: an arrow has, in addition, an identity-on-objects functor
from a cartesian base and a $\first$ operator. We take the tensor to be the cartesian product
$\times$ on $\Cont$. For a pure morphism $\varphi\colon p\to q$ set
$\arr(\varphi)=\eta^T_q\circ\varphi\circ\varepsilon_p\colon G_M p\to T_M q$.

\begin{proposition}[The pure functor]\label{prop:arr}
$\arr$ is an identity-on-objects functor $\Cont\to\Arr_M$: $\arr(\id)$ is the biKleisli
identity, and $\arr(\psi\circ\varphi)=\arr(\varphi)\,\rcomp\,\arr(\psi)$.
Its image is the pure/central subcategory.
\end{proposition}
\begin{proof}
Functoriality is the counit/unit calculation: expanding the composite and using
$\varepsilon\circ\delta=\id$ together with the entwining triangles E1$'$
($\kappa\circ G_M\eta^T=\eta^T G_M$) and E3$'$ ($T_M\varepsilon\circ\kappa=\varepsilon T_M$),
both of which hold for \emph{every} $M$ (\S\ref{sec:onedirection}), collapses the middle
$\eta^T\cdots\varepsilon$ to give $\eta^T\circ\psi\circ\varphi\circ\varepsilon
=\arr(\psi\circ\varphi)$.
\end{proof}

The coeffect side is unconditionally well-behaved; the effect side is where branching bites
again.

\begin{lemma}[Coeffects are always costrong]\label{lem:costrong}
For every $M$ there is a natural costrength $\sigma_{p,c}\colon G_M(p\times c)\to G_M p\times
c$, forward $\id$, backward at $(v,t)$ the map $M(Qv)\sqcup Q_c t\to M(Qv\sqcup Q_c t)$ that is
$M(\mathrm{inl})$ on the coeffect summand and $\eta^M\circ\mathrm{inr}$ on the fresh summand.
\end{lemma}
\begin{proof}
$G_M$-positions are a \emph{single} $M$-structure $M(Ps)$, with no product over leaves, so the
map is $M$ applied to $\mathrm{inl}$ and the unit on the new $c$-factor; no branching arises.
Naturality is functoriality of $M$ and naturality of $\eta^M$.
\end{proof}

\begin{lemma}[Effects are strong iff non-branching]\label{lem:strong}
The effect monad $T_M$ admits a \emph{natural} tensorial strength $\tau_{q,c}\colon T_M
q\times c\to T_M(q\times c)$ for $\times$ if and only if $M$ is non-branching.
\end{lemma}
\begin{proof}
Backward, $\tau$ is the distributivity map
$d_{m,t}\colon\prod_{b\in\lv m}\bigl(Q(v_b)\sqcup Q_c t\bigr)\to\bigl(\prod_b Q(v_b)\bigr)
\sqcup Q_c t$.

$(\Leftarrow)$ If every $m$ has $|\lv m|\le 1$ then this is, at a nullary shape, the coprojection
$1\to 1\sqcup Q_c t$, and at a unary shape the identity on $Q(v_0)\sqcup Q_c t$ (a length-$1$
product is its unique factor). No choice is made, so $\tau$ is natural and satisfies the
strength unit and multiplication axioms.

$(\Rightarrow)$ Suppose some $m_0$ has $n=|\lv m_0|\ge 2$. The family $d$ must be natural in
$q$ and $c$, i.e.\ a natural transformation of the functors $L(A,C)=\prod_b(A_b\sqcup C)$ and
$R(A,C)=(\prod_b A_b)\sqcup C$ of the variables $A_b:=Q(v_b)$, $C:=Q_c t$, and equivariant
under the reindexings of $\lv m_0$ realised by $M$'s functoriality. Specialise every
$A_b=\varnothing$: then $L(\varnothing,C)=C^{\,n}$ and $R(\varnothing,C)=C$, so $d$ restricts
to a natural map $C^{\,n}\Rightarrow C$. Since $C^{\,n}=\Hom(n,-)$, Yoneda gives
$\Nat(\Hom(n,-),\mathrm{Id})\cong n$: the natural maps $C^{\,n}\to C$ are exactly the $n$ leaf
projections, so $d|_{A=\varnothing}=\pi_i$ for a \emph{fixed} distinguished leaf $i$. That
distinguished leaf is the contradiction. For \emph{symmetric} branching (e.g.\ $\Pf$),
$M$'s functoriality realises the transposition of two leaves of $m_0$ as an automorphism
fixing the shape; equivariance forces $\pi_i=\pi_j$, impossible on $C^{\,n}$ with $|C|\ge2$.
For \emph{ordered} branching (e.g.\ $\List$), a leaf-reindexing morphism breaks the fixed
choice likewise. Hence no natural $d$, so no natural $\tau$.
\end{proof}

\begin{remark}[Two faces of branching]
A \emph{total} strength does exist for branching $M$ --- the ``priority'' rule (all-$\mathrm{
inl}$ to the product, else the leftmost $\mathrm{inr}$ value) is total and even satisfies the
strength unit and multiplication axioms --- but it is \emph{not natural}: it hard-codes the
distinguished leaf that naturality forbids. So the strength obstruction is genuinely distinct
from the associativity obstruction of Theorem~\ref{thm:arrowsA}: the latter is $\mu^T$
\emph{merging} (E2$'$, union-of-products), the former is leaf-\emph{symmetry} (naturality).
Branching disables the arrow through two independent axioms, both equivalent to arity $\le1$.
\end{remark}

\begin{theorem}[The arrow is Freyd iff non-branching]\label{thm:arrowsB}
For non-branching $M$, with $\sigma_G$ (Lemma~\ref{lem:costrong}) and $\tau_T$
(Lemma~\ref{lem:strong}) and $\first(f)=\tau_{q,c}\circ(f\times\id_c)\circ\sigma_{p,c}$, the
data $(\Arr_M,\arr,\rcomp,\first)$ satisfy all the Hughes arrow laws;
$\Arr_M$ is a Hughes arrow, equivalently a Freyd category over the cartesian base
$(\Cont,\times)$. For branching $M$, $\Arr_M$ is not even a category
(Theorem~\ref{thm:arrowsA}), and $\first$ is not definable ($T_M$ has no natural strength,
Lemma~\ref{lem:strong}). Hence
\[
  \boxed{\ \Arr_M\ \text{is a Hughes arrow / Freyd category}\iff M\ \text{is non-branching}.\ }
\]
\end{theorem}
\begin{proof}
The constructions $\arr,\first$ are uniform in $M$; the abstract packaging is the standard
theorem that the biKleisli category of a strong monad and a costrong comonad with a compatible
distributive law is a Freyd/arrow category \cite{UustaluVene08,PowerRobinson97,Atkey11,
JacobsHeunenHasuo09}. The container-specific content --- that $\sigma_G,\tau_T,\kappa$ satisfy
the required (co)strength and compatibility coherences --- holds for non-branching $M$ by
Lemmas~\ref{lem:costrong}--\ref{lem:strong} and Theorem~\ref{thm:arrowsA}, and is verified
exhaustively for the representatives $\Maybe$ and $\Writer/\mathbb{Z}_2$, which span the
normal form of \S\ref{sec:classify}. The branching direction is Theorem~\ref{thm:arrowsA} and
Lemma~\ref{lem:strong}.
\end{proof}

\begin{remark}[$\Arr_M$ is Atkey's biKleisli Arrow; branching is transverse to Atkey's index]\label{rem:atkey}
Atkey \cite{Atkey11} shows that Hughes arrows are \emph{not} the same thing as Freyd
categories: an arrow is strictly more general --- it admits a second, comonad-structured
input that a Freyd category lacks --- and the precise correspondence is \emph{arrows $\cong$
closed \textbf{indexed} Freyd categories}, the extra generality controlled by an \emph{index}
that records the comonadic input. His running construction (\S2, after Uustalu--Vene
\cite{UustaluVene08}) is exactly ours: from a comonad $W$, a monad $T$, and a distributive
law $WT\Rightarrow TW$ he forms the biKleisli arrow $\mathrm{Ar}(x,y)=[Wx\to Ty]$; with
$W=G_M$, $T=T_M$, and $\lambda=\kappa$ this is \emph{verbatim} our
$\Arr_M(p,q)=\Cont(G_M p,\,T_M q)$.

It is tempting to read Theorem~\ref{thm:arrowsB} as an \emph{index-collapse} --- non-branching
$M$ trivialising Atkey's index so that $\Arr_M$ becomes a plain (non-indexed) Freyd category
--- but that reading is \emph{false}, and the way it fails is informative. Atkey's index
records the \emph{coeffect} input $W=G_M$, and $G_M$ is trivial exactly when $M$ is:
$G_M=\mathrm{Id}$ iff $M=\mathrm{Id}$, immediate from $G_M\cont S P=\cont S{M\circ P}$.
Index-collapse is therefore the single condition $M=\mathrm{Id}$, \emph{strictly stronger}
than non-branching: $\Maybe$ and $\Writer$ are non-branching with $G_M\neq\mathrm{Id}$, and
Theorem~\ref{thm:arrowsB} still makes their $\Arr_M$ a genuine Freyd category --- through the
\emph{strength} of $T_M$, not through any collapse of the coeffect index. The branching axis
and Atkey's index axis are thus \emph{orthogonal}: index-collapse ($M=\mathrm{Id}$) is a proper
sub-condition of non-branching, and Atkey's index measures a coeffect that persists across the
entire non-branching class. \emph{Non-branching is not index-collapse.}

What branching controls is whether the biKleisli construction lands in Atkey's landscape at
all: for branching $M$ the compositor $\kappa$ is not even a distributive law (E2$'$ fails,
Theorem~\ref{thm:arrowsA}), so $\Arr_M$ degrades below an indexed Freyd category --- below even
an arrow. The branching obstruction sits on the \emph{arity} axis, transverse to Atkey's
monad~$\subset$~Freyd~$\subset$~indexed-Freyd~$\subset$~arrow stratification of the coeffect;
whether the resulting (Boolean) arity dichotomy admits any graded refinement is taken up in
\S\ref{sec:conclusion}.
\end{remark}

\begin{remark}[The pure and effectful sub-calculi]
The image of $\arr$ is the pure/central subcategory. The \emph{effect} sub-calculus is the
Kleisli category $\Kl(T_M)$ (collapse $G_M$ by $\varepsilon$): container programs with the
Ahman--Bauer effect. The \emph{coeffect} sub-calculus is the coKleisli category $\coKl(G_M)$
(collapse $T_M$ by $\eta^T$): container programs reading the transfer-comonad context. The
arrow $\Arr_M$ \emph{fuses} them, and Theorem~\ref{thm:arrowsA} says the fusion is coherent
exactly when $M$ does not branch.
\end{remark}

%============================================================================
\section{The positive class: writer with absorbing exceptions}\label{sec:classify}
%============================================================================

Theorems~\ref{thm:arrowsA}--\ref{thm:arrowsB} leave a name to be attached: which monads are
the non-branching ones? We classify them completely.

\begin{theorem}[Classification]\label{thm:classify}
For a cartesian $\prod$-cointerpretation $\Set$-monad $M$ the following are equivalent:
\begin{enumerate}
  \item the arrow category $\Arr_M$ exists;
  \item $M$ is non-branching;
  \item $M\cong E + A\times(-)$ as a functor, for some sets $E,A$;
  \item $M$ is a \emph{writer-with-absorbing-exceptions} monad: $MX = E + A\times X$ with
  $(A,\cdot,e)$ a monoid, $(E,\odot)$ a left $A$-set, writer unit $\eta_X(x)=(e,x)$, and
  \[
    \mu_X\colon\ \mathrm{inl}\,e'\mapsto\mathrm{inl}\,e',\quad
    \mathrm{inr}(a,\mathrm{inl}\,e')\mapsto\mathrm{inl}(a\odot e'),\quad
    \mathrm{inr}(a,\mathrm{inr}(a',x))\mapsto\mathrm{inr}(a\cdot a',x).
  \]
\end{enumerate}
Moreover, for every such $M$ and every container $p$ the compositor $\kappa$ satisfies all
four axioms E1$'$--E4$'$; in particular the associativity axiom E2$'$ holds throughout the
class, so $\Arr_M$ is a category for the entire class, not only for $\Maybe$ and $\Writer$.
\end{theorem}

\begin{proof}
$(1)\Leftrightarrow(2)$ is Theorem~\ref{thm:arrowsA}.

$(2)\Leftrightarrow(3)$: a polynomial functor is $MX=\sum_{s\in M1}X^{\mathrm{ar}(s)}$. If
$|\mathrm{ar}(s)|\le1$ for all $s$, partition the shapes as $E=\{s:\mathrm{ar}(s)=
\varnothing\}$ and $A=\{s:|\mathrm{ar}(s)|=1\}$; then $X^{\mathrm{ar}(s)}$ is $1$ on $E$ and
$X$ on $A$, so $MX\cong E + A\times X$ naturally. Conversely $E+A\times(-)$ has all arities
$\le1$.

$(3)\Rightarrow(4)$: let $MX=E+A\times X$ carry a monad. By naturality (Yoneda for the
identity domain), the unit $\eta$ is either a constant $x\mapsto\mathrm{inl}\,e_0$ or
$x\mapsto\mathrm{inr}(a_0,x)$; the constant form violates the left-unit law, so
$\eta_X(x)=\mathrm{inr}(e_0,x)$ for a distinguished $e_0\in A$ (the writer unit). Naturality
similarly forces $\mu$ to the schematic form above, carrying the variable $x$ exactly when the
outer and inner $A$-labels multiply into $A$. Writing $N:=E\sqcup A$ with the induced binary
operation, the unit laws make $e_0$ a two-sided unit in $A$, and the associativity of $\mu$
unwinds to the associativity of $N$; $E$ is automatically a two-sided ideal of left zeros. In
the \emph{cartesian} case $\mu^M$ preserves the leaf-count, which rules out \emph{aborting}
($a\cdot a'\in E$ for $a,a'\in A$ would destroy the single leaf), so $A$ is a submonoid and
$\sigma$ restricts to a unital, associative action $\odot\colon A\times E\to E$: $E$ is a left
$A$-set, statement (4). $(4)\Rightarrow(3)$ is immediate.

Finally, the last clause. Non-branching means every $P^\star$-product occurring anywhere is a
product over $\le1$ factor, so $\kappa$'s backward map is, shape by shape, either the identity
(one leaf: rewrap the single factor) or the monad unit $\eta^M\colon1\to M1$ (no leaf: the
empty product). Reading each axiom E1$'$--E4$'$ as an equality of $\Cont$-morphisms: on shapes
all four reduce to the monad/comonad shape-laws of $M$; on positions the only nontrivial
content sits at unary shapes, where $\kappa=\id$ and the axiom collapses to the corresponding
structural law of $M$ on the single leaf --- for E2$'$, the associativity of $\cdot$ in $A$,
which holds. The $\ge2$-leaf comparison that breaks E2$'$ for branching $M$ never forms.
\end{proof}

\begin{remark}[The Set-monad level, and why aborting is excluded]
One level up, the functors $E+A\times(-)$ that carry \emph{any} monad structure are exactly
the monoids $N=E\sqcup A$ with unit in $A$ and $E$ a two-sided ideal of left zeros --- the
multiplication of two unary shapes may \emph{abort} into $E$ (e.g.\ the nilpotent
$\{e,z,0\}$ with $z^2=0\in E$, giving $MX=1+2\times X$). These aborting monads are genuine
$\Set$-monads but are \emph{not cartesian}: their $\mu$ destroys a leaf, so Ahman--Bauer's
$T_M$ is undefined for them and they fall outside the arrow story. Within the class where
$T_M$ exists, the non-branching monads are exactly the non-aborting ones --- the
writer-with-absorbing-exceptions monads. (With trivial action $\odot$, $E+A\times X$ is the
writer transformer over the exception monad, $\mathsf{WriterT}_A(\mathsf{Exc}_E)(X)$; general
$\odot$ twists the discard by a left $A$-action on $E$.)
\end{remark}

\begin{remark}[Scope]\label{rem:scope}
Both $\lambda$ and $\kappa$, and the strength $\tau$, use the product structure of $P^\star$;
whether a canonical law exists for a general (non-$\prod$) weak Mendler algebra is open, and
out of scope. The Freyd base is the cartesian $\times$; a monoidal (non-cartesian) arrow over
the Dirichlet tensor $\otimes$ would need a different strength story and is a separate
question. No finite $\prod$-cointerpretation monad is simultaneously branching and
non-commutative ($\List$ is the witness, but infinite), so that corner of the entwining
theorem is proved by naturality and untested computationally.
\end{remark}

\subsection{Related conditions: non-branching is genuinely new}\label{sec:related-conditions}
The dividing line of Theorems~\ref{thm:arrowsA}--\ref{thm:classify} is \emph{non-branching}.
One should check that this is not a disguised form of one of the two classical properties of a
$\Set$-monad --- Kock \emph{commutativity} \cite{Kock70,Kock71} and Kock/Jacobs
\emph{affineness} $M1\cong1$ \cite{Jacobs16} (the sense of Remark~\ref{rem:affine-clash}, not
the polynomial-degree sense we deliberately avoid). It is not: writing
\[
  \text{P1}=\text{non-branching},\qquad \text{P2}=\text{commutative},\qquad
  \text{P3}=\text{affine }(M1\cong1),
\]
the three are pairwise logically independent, with a single implication among them pointing the
other way (P1$\wedge$P3$\Rightarrow$P2). The commutativity of the members of the positive class
is itself sharp:

\begin{lemma}[Commutativity criterion for the class]\label{lem:commcrit}
A writer-with-absorbing-exceptions monad $MX=E+A\times X$ (Theorem~\ref{thm:classify}) is
commutative if and only if \textup{(i)} $A$ is a commutative monoid, \textup{(ii)} $|E|\le1$,
and \textup{(iii)} the action $\odot$ is trivial. In particular $\Maybe$ ($|E|=1$, $A=1$) and
$\Writer_A$ with $A$ abelian are commutative, whereas $\Writer_{N_3}$ over a non-commutative
monoid $N_3$ and the exception monad $E+(-)$ with $|E|\ge2$ are not.
\end{lemma}
\begin{proof}
$M$ is commutative iff Kock's two canonical maps $MX\times MY\to M(X\times Y)$ agree.
Evaluating both on the four kinds of pair in $(E+A\times X)\times(E+A\times Y)$: on
$(\mathrm{inl}\,e,\mathrm{inl}\,e')$ they return $\mathrm{inl}\,e$ and $\mathrm{inl}\,e'$
(equal for all inputs iff $|E|\le1$); on $(\mathrm{inl}\,e,\mathrm{inr}(b,y))$ they return
$\mathrm{inl}\,e$ and $\mathrm{inl}(b\odot e)$, and symmetrically on the mirror pair (equal iff
$\odot$ is trivial); on $(\mathrm{inr}(a,x),\mathrm{inr}(b,y))$ they return
$\mathrm{inr}(a\cdot b,(x,y))$ and $\mathrm{inr}(b\cdot a,(x,y))$ (equal iff $A$ is
commutative). All four agree iff (i)--(iii).
\end{proof}

\begin{proposition}[Non-branching is independent of commutativity and affineness]\label{prop:independence}
Each of the three $2\times2$ faces of the cube $(\mathrm{P1},\mathrm{P2},\mathrm{P3})$ is fully
realised by explicit $\Set$-monads; hence no one of the three properties implies or precludes
another. Concretely:
\[
  \renewcommand{\arraystretch}{1.3}
  \begin{array}{c|c|l}
    (\mathrm{P1},\mathrm{P2},\mathrm{P3}) & \text{witness } M & \text{values}\\\hline
    (T,T,T) & \mathrm{Id} & \text{ar}\le1,\ \text{comm},\ M1=1\\
    (T,T,F) & \Maybe=1+(-) & \text{ar}\le1,\ \text{comm},\ M1=2\\
    (T,F,T) & \text{--- impossible ---} & \text{Remark~\ref{rem:theoremC}}\\
    (T,F,F) & \Writer_{N_3}=N_3\times(-) & \text{ar}\le1,\ \text{non-comm},\ M1=3\\
    (F,T,T) & \mathcal P^{+}\ (\text{non-empty powerset}) & \text{branch},\ \text{comm},\ M1=1\\
    (F,T,F) & \Pf & \text{branch},\ \text{comm},\ M1=2\\
    (F,F,T) & \text{free idempotent magma} & \text{branch},\ \text{non-comm},\ M1=1\\
    (F,F,F) & \text{free magma (binary trees)} & \text{branch},\ \text{non-comm},\ M1=\infty
  \end{array}
\]
Here $N_3=\{1,a,b\}$ is the three-element non-commutative monoid with $a\cdot x=a$ and
$b\cdot x=b$ for $x\in\{a,b\}$.
\end{proposition}
\begin{proof}
Each $2\times2$ face is a projection of the table (e.g.\ P1$\times$P2 is realised by
$\mathrm{Id},\Writer_{N_3},\Pf$, and the free magma), so every pair of properties attains all
four truth-values and pairwise independence follows. The property values are established as:
\textbf{P1} by the support certificate of \S\ref{sec:prelim-lifting} (the $E+A\times(-)$
witnesses have all supports $\le1$; $\mathcal P^{+},\Pf$ and both magmas carry a support-$2$
element); \textbf{P2} by Lemma~\ref{lem:commcrit} for the non-branching witnesses, by the
classical commutativity of the semilattice monads $\mathcal P^{+},\Pf$ \cite{Kock71}, and by
the failure of the medial/interchange law $(a*b)*(c*d)=(a*c)*(b*d)$ --- forced by Kock's
theorem that a commutative monad has homomorphic operations --- in explicit small models of the
two magmas; \textbf{P3} by the displayed values of $|M1|$. Full witnesses, the four-row
commutativity chase, and an exhaustive machine sweep are in the author's proof notes.
\end{proof}

\begin{remark}[The one dependency]\label{rem:theoremC}
The three properties are \emph{not} jointly independent: the cube has exactly one empty cell,
$(\mathrm{P1},\lnot\mathrm{P2},\mathrm{P3})$, because \textbf{non-branching and affine imply
commutative}. If $M\cong E+A\times(-)$ is non-branching then $M1\cong E+A$, so affineness
$|M1|=1$ forces $(|E|,|A|)\in\{(0,1),(1,0)\}$, i.e.\ $M\cong\mathrm{Id}$ or the constant monad
$1$, both commutative. Equivalently, \emph{every non-commutative affine monad is branching}:
affineness collapses the non-branching world to the two trivial monads and so cannot express
non-branching. This lone implication --- pointing away from non-branching --- is the only
logical relation among P1, P2, P3, which confirms that ``non-branching'' is a genuinely new
dividing line rather than a restatement of a classical monad condition.
\end{remark}

%============================================================================
\section{A machine-checked instance}\label{sec:lean}
%============================================================================

The non-branching monad $M=\Maybe$ furnishes the smallest genuinely effectful arrow category,
and it has been formalised in Lean~4. The development \texttt{BiKleisliMaybe.lean} constructs
the effect monad $T_{\Maybe}\cont S P=(1+S,\ P^\star)$ as an instance of an abstract
\texttt{MixedDistrib} structure --- monad $T$, comonad $G$, compositor $\kappa$, and the four
axioms E1$'$--E4$'$ including the associativity axiom E2$'$ --- and derives the biKleisli
composition together with its associativity law \texttt{acomp\_assoc} from those axioms, with
no \texttt{sorry}. This is, to the author's knowledge, the first machine-checked
\emph{associative} effect--coeffect arrow category on containers for a concrete non-branching
monad. It verifies the load-bearing content of Theorem~\ref{thm:arrowsA} in the smallest
non-trivial case; the general E2$'$ across the whole class (Theorem~\ref{thm:classify}) is
proved in coordinates but not yet Lean-certified, and remains a natural next target.

%============================================================================
\section{Related work}\label{sec:related}
%============================================================================

\paragraph{The two feeds, and their sources.}
The effect monad $T_M$ is Ahman--Bauer \cite[Thm.~6.3]{AhmanBauer24}; the coeffect comonad
$G_M$ is the transfer (the author's, proved and Lean-verified). Neither paper forms the
composites $T_M G_M$ / $G_M T_M$ or a distributive law between the feeds. The entwining and
arrow formalism is classical: Beck \cite{Beck69}; Power--Watanabe \cite{PowerWatanabe02};
Brzezi\'nski--Majid \cite{BrzezinskiMajid98}; Uustalu--Vene \cite{UustaluVene08}; Hughes
\cite{Hughes00}; Power--Robinson \cite{PowerRobinson97}; Atkey \cite{Atkey11};
Jacobs--Heunen--Hasuo \cite{JacobsHeunenHasuo09}; Turi--Plotkin \cite{TuriPlotkin97}. We
deploy it where the base monad sits on positions and branching is the decisive parameter.

\paragraph{Distributive laws of container (co)monads (Purdy--Damato).}
Purdy and Damato \cite{PurdyDamato25} characterise distributive laws between
container-monads and container-comonads \emph{on $\Set$}, including mixed
monadic-over-directed laws, by an explicit data presentation. By contrast our $T_M$ and $G_M$
are a monad and a comonad on the category $\Cont$ itself, obtained from a single $\Set$-monad
$M$ by transferring it to shapes and to positions respectively; their (co)monads are container
functors on $\Set$, whereas ours are endofunctors of $\Cont$, one level up. Consequently our
mixed law $\lambda$ and its arrow-orientation $\kappa$ live in $[\Cont,\Cont]$ and are not
instances of their classification. (Their Example~22 links monadic-container laws to matching
pairs of monoid actions in the sense of Zappa--Sz\'ep, adjacent to the directed mode discussed
below but at the level of monoids rather than of the entwining.)

\paragraph{Interaction laws (Katsumata--Rivas--Uustalu).}
The nearest neighbour by intent is the theory of \emph{interaction laws of monads and comonads}
\cite{KRU20}, which answers effect/coeffect interaction by a \emph{pairing}
$\psi\colon TX\times DY\to X\times Y$, realised as a monoid object in a Chu construction over
$([\mathcal C,\mathcal C],\mathrm{Day})$. That is emphatically \emph{not} a distributive law
--- distributive laws are \emph{inputs} to their machinery, never outputs --- so the two
accounts are orthogonal in kind, a pairing versus a compositor. What they share is only the
\emph{location} of the obstruction: their degeneracy theorems locate the same branching wall
from the side of extensive categories that our arity criterion locates from the side of
polynomial functors --- convergent evidence that branching is the universal obstruction to
composing effects with coeffects.

\paragraph{Neighbours read only in outline.}
Two further lines of work are close in intent and deserve a detailed comparison we defer here
to a later reading. Modern higher-order \emph{bialgebraic denotational semantics} (Goncharov,
Milius, Tsampas, Urbat, and collaborators, ICFP 2026) builds bialgebras from a syntax
$\Sigma$-algebra and a behaviour $B$-coalgebra glued by higher-order GSOS laws via locally
final coalgebras --- a different object from our monad-over-comonad mixed distributive law, in
which both $T_M$ and $G_M$ arise as the two feeds of one $\Set$-monad rather than as separate
syntax and behaviour functors. Separately, Dumas, Duval and Reynaud (\emph{Patterns for
computational effects arising from a monad or a comonad}, 2013) pair a monad with its
associated comonad as two dual but \emph{separate} equational proof-patterns (Kleisli versus
coKleisli), whereas we entwine an independent shapes-monad and positions-comonad via a genuine
distributive law. We flag both as neighbours rather than cite them as sources, having so far
consulted them only in outline.

\paragraph{Three modes of composition.}
The effect--coeffect arrow is one of three container-native modes of welding compositional
systems the author has studied, and it is the branching-obstructed one. The other two are the
Zappa--Sz\'ep product of two directed containers (obstructed by a cohomology class
$[\omega]\in H^2$) and the state-graded category of stateful ``Workers'' (unobstructed, the
grade accumulating). The three obstructions --- a cohomology class, nothing, and a Boolean
arity condition --- are three genuinely different kinds of mathematics for three ways of
gluing; they do not unify into a master theorem, and the honest synthesis is that
\emph{whether} two systems compose is, in each mode, a computable invariant of the parts,
though which invariant depends on the axis.

%============================================================================
\section{Conclusion}\label{sec:conclusion}
%============================================================================

One monad $M$ on $\Set$ feeds a container along two axes: forward on shapes it is the effect
monad $T_M$, backward on positions it is the coeffect comonad $G_M$. We have shown that these
two feeds always entwine --- the oplax product-comparison of $M$, read on positions, is a
mixed distributive law $\lambda\colon T_M G_M\Rightarrow G_M T_M$ for every $M$ --- giving a
Turi--Plotkin $\lambda$-bialgebra in which $G_M$ acts on $T_M$-algebras and $T_M$ on
$G_M$-coalgebras, nondeterminism and all. The reverse compositor $\kappa$ that would let one
\emph{sequentially compose} effect--coeffect programs $G_M p\to T_M q$ into a Hughes arrow /
Freyd category exists if and only if $M$ is non-branching, and that class is exactly the
writer-with-absorbing-exceptions monads $E+A\times(-)$. The same monad, the same two feeds:
the unification of effects and coeffects is unconditional as a bialgebra and conditional as an
arrow, and branching is precisely what separates the two --- obstructing the orientation of the
interaction, not the interaction itself.

Several questions remain open and are, to us, the natural continuations.
\begin{itemize}
  \item \emph{Beyond the $\prod$-cointerpretation.} Both laws use the product structure of
  $P^\star$. Is there a canonical $\lambda$ for a general (non-$\prod$) weak Mendler algebra?
  There is no evident map $M(P^\star)\to(MP)^\star$ without the product, so we suspect not, but
  a proof either way is missing (Remark~\ref{rem:scope}).
  \item \emph{A graded refinement?} Read through Atkey's arrows~$\cong$~indexed Freyd
  categories (Remark~\ref{rem:atkey}), one might hope the Boolean dichotomy ``$\Arr_M$ is a
  category iff arity ${\le}1$'' is the bottom rung of a \emph{graded} tower indexed by arity.
  A companion analysis rules out the two natural gradings. Along the effect axis the arity
  invariant is already two-valued --- a cartesian $\Set$-monad has maximal arity in
  $\{{\le}1,\infty\}$, with no finite rung between, since an arity-$n$ shape substituted into
  itself has arity $n^2$ --- so ``arity ${\le}n$'' names no monads for $2\le n<\infty$; and the
  leaf-count grade on arrows is destroyed by the very merging in $\mu^{T_M}$ that branching
  forces. A grade that survives, if one exists, must live on the \emph{coeffect}: Atkey's index
  measures $G_M$, nontrivial for every $M\neq\mathrm{Id}$ and independent of branching, and
  whether a graded comonad on $G_M$ organises $\Arr_M$ along that axis is left open.
  \item \emph{The Dirichlet arrow.} The Freyd reading fixes the cartesian $\times$. Does the
  Dirichlet tensor $\otimes$ on $\Cont$ carry a \emph{monoidal} (non-cartesian) effect--coeffect
  arrow, with a different strength story? The costrength of $G_M$ is the easy half; the effect
  strength for $\otimes$ is the open one.
  \item \emph{Full formalisation.} The Maybe arrow category is machine-checked
  (\S\ref{sec:lean}); the associativity axiom E2$'$ across the whole writer-with-exceptions
  class (Theorem~\ref{thm:classify}) is proved in coordinates but not yet in Lean. Certifying
  it, and the general mult-$T$ index-chase for $\lambda$ (Remark~\ref{rem:e2gap}), would close
  the two mechanical gaps.
\end{itemize}

%============================================================================
\begin{thebibliography}{99}

\bibitem{AhmanBauer24} D.~Ahman, A.~Bauer.
\emph{Comodule representations of second-order functionals}.
\href{https://arxiv.org/abs/2409.17664}{arXiv:2409.17664};
J.\ Log.\ Algebraic Methods Program.\ \textbf{146} (2025).

\bibitem{Atkey11} R.~Atkey.
\emph{What is a categorical model of arrows?}
Electronic Notes in Theoretical Computer Science \textbf{229}(5) (2011) 19--37.

\bibitem{PurdyDamato25} C.~Purdy, D.~Damato.
\emph{Distributive laws of monadic containers}.
\href{https://arxiv.org/abs/2503.17191}{arXiv:2503.17191}; Proc.\ CALCO 2025.

\bibitem{Beck69} J.~Beck.
\emph{Distributive laws}.
In: Seminar on Triples and Categorical Homology Theory, Lecture Notes in Mathematics
\textbf{80}, Springer (1969) 119--140.

\bibitem{BrzezinskiMajid98} T.~Brzezi\'nski, S.~Majid.
\emph{Coalgebra bundles}.
Communications in Mathematical Physics \textbf{191} (1998) 467--492.

\bibitem{Hughes00} J.~Hughes.
\emph{Generalising monads to arrows}.
Science of Computer Programming \textbf{37}(1--3) (2000) 67--111.

\bibitem{Jacobs16} B.~Jacobs.
\emph{Affine monads and side-effect-freeness}.
Proc.\ CMCS 2016, LNCS \textbf{9608}, Springer (2016) 53--72.

\bibitem{JacobsHeunenHasuo09} B.~Jacobs, C.~Heunen, I.~Hasuo.
\emph{Categorical semantics for arrows}.
Journal of Functional Programming \textbf{19}(3--4) (2009) 403--438.

\bibitem{KRU20} S.~Katsumata, E.~Rivas, T.~Uustalu.
\emph{Interaction laws of monads and comonads}.
Proc.\ FoSSaCS 2020, LNCS \textbf{12077}, 356--375;
\href{https://arxiv.org/abs/1912.13477}{arXiv:1912.13477}.

\bibitem{Kock70} A.~Kock.
\emph{Monads on symmetric monoidal closed categories}.
Archiv der Mathematik \textbf{21} (1970) 1--10.

\bibitem{Kock71} A.~Kock.
\emph{Closed categories generated by commutative monads}.
Journal of the Australian Mathematical Society \textbf{12}(4) (1971) 405--424.

\bibitem{PowerRobinson97} J.~Power, E.~Robinson.
\emph{Premonoidal categories and notions of computation}.
Mathematical Structures in Computer Science \textbf{7}(5) (1997) 453--468.

\bibitem{PowerWatanabe02} J.~Power, H.~Watanabe.
\emph{Combining a monad and a comonad}.
Theoretical Computer Science \textbf{280} (2002) 137--162.

\bibitem{TuriPlotkin97} D.~Turi, G.~Plotkin.
\emph{Towards a mathematical operational semantics}.
Proc.\ LICS 1997, IEEE (1997) 280--291.

\bibitem{UustaluVene08} T.~Uustalu, V.~Vene.
\emph{Comonadic notions of computation}.
Electronic Notes in Theoretical Computer Science \textbf{203}(5) (2008) 263--284.

\end{thebibliography}

\end{document}
